\begin{afterword}
\summaryhead

The post starts from the purpose we see in nature: a rabbit's legs built for running, a fox's jaws for tearing. Before Darwin, design by God was a reasonable explanation. But nature as a whole fits one benevolent designer poorly. Foxes are designed to catch rabbits and rabbits to escape foxes, which suggests competing gods; and much design is cruel, like the ichneumon wasp whose larvae eat caterpillars alive.

Evolution explains both, the author argues, if it is understood. There is no ``Evolution Fairy'' deciding which genes would be helpful. A gene becomes more common when its effects cause more copies of it to reach the next generation, and this happens separately in each population. So nature is at war with itself, cruelty that would cost nothing to remove persists, and a design like the backward human retina cannot be fixed one mutation at a time. Evolution is closer to a god than to random chance, and godlike in some ways, but it is mindless, divided and not nice: Lovecraft's Azathoth. Religionists who claim to believe in a vague deity, or to await science's discovery of God, should have recognized it.

\responsehead

The biology is the standard gene-centred account, stated clearly and with its hedges. The post keeps to cause and effect, marks its simplifications, and gets the arithmetic of kin selection right. Its line that evolution ``is closer to God than it is to pure random entropy'' is a correct and useful way of saying that selection, unlike mutation, is not random.

The theological argument is older than the post suggests. Twice the post wonders whether anyone had noticed what nature's conflicts and cruelty imply, and the first time it answers ``Probably not.'' In Hume's \textsc{Dialogues Concerning Natural Religion} (1779), Philo asks why several deities might not have framed the world, as many workers build a ship. Later, looking at how ``hostile and destructive to each other'' living things are, he says that the Manichaean system of a good and an evil power ``has more probability than the common hypothesis,'' and then prefers first causes with ``neither goodness nor malice.'' That is close to the post's Azathoth. What the post adds is the mechanism, gene frequencies, which Hume did not have. The step from nature's conflict and cruelty to an indifferent maker was not new.

The principle that a theory's strength lies in what it forbids is Popper's, uncredited here as in earlier posts.

The ending sits uneasily with the post's own account. It charges religionists who claim a vague deity with failing to recognize evolution as their creator. None is named or quoted. The charge also needs a notion of deity that a mindless process can meet, and the post itself says that ``Evolution is not a God'' and that ``the Maker has no mind.'' A believer who declined to call evolution a god was agreeing with the post. The same holds for the religionists said to be waiting for science to discover God.

In short: a clear account of selection without foresight, set inside a theological argument whose main moves Hume had made in 1779, and closed with a charge against unnamed religionists that the post's own ``Evolution is not a God'' answers.
\end{afterword}
